\section{Conclusion}

This paper introduced a finite connected intervention calculus for
converter-rich power-system DAEs. The defining choice is physical
re-equilibration: each subset of a finite action portfolio receives its own ac
solution, dynamic initialization, and descriptor operator. Direct M\"obius
differences isolate pair and triple externalities, and connected total
derivatives reconstruct them over the full action hypercube.

The registered IEEE 39-bus campaign supports five bounded conclusions.
Finite pair and genuine third-order externalities exist reproducibly. The
connected reconstruction closes all locked-holdout cases within five
conservative uncertainty budgets. Exact Schur anatomy separates algebraic,
mass-scaling, and finite-pole contributions to roundoff. Reproducible
finite-pole externalities exist, and certified contours localize pole motion
for three retained coalitions. Re-equilibration can be quantitatively
decisive: the frozen-affine approximation underestimates A7--A8 by about
$6.3\times$ at the OP00 reference point.

The experiment also establishes a firm limit. None of the selected
externalities met the preregistered nominal or load-conditioned damping
materiality criteria. The weak-grid test was baseline-ineligible and remains
unresolved. A separate audit showed why: the externality-dominant modes are
well-damped converter-control families, whereas the actual limiting mode is a
roughly 1.22-Hz synchronous-electromechanical family with only microunit-scale
finite damping composition.

The resulting conclusion is deliberately narrower than an interaction-based
stability claim: finite dynamic externalities can be real, higher order, and
reproducible without materially consuming the system's registered modal
margin. Detecting interaction structure is therefore a useful diagnostic step
when finite portfolio composition is the decision question, but it is not a
substitute for identifying
the constraint-setting mode or for demonstrating a mechanism-specific
intervention.
